\documentclass[11pt]{article}

% ============================================================
% Packages
% ============================================================
\usepackage[margin=1in]{geometry}
\usepackage{amsmath,amssymb,mathtools,amsthm}
\usepackage{enumitem}
\usepackage{hyperref}
\usepackage{xcolor}
\usepackage{microtype}
\usepackage{booktabs}
\usepackage{cleveref}

% ============================================================
% Hyperref setup
% ============================================================
\hypersetup{
  colorlinks=true,
  linkcolor=blue!60!black,
  citecolor=blue!60!black,
  urlcolor=blue!60!black
}

% ============================================================
% Macros
% ============================================================
\newcommand{\R}{\mathbb{R}}
\newcommand{\N}{\mathbb{N}}
\newcommand{\F}{\mathcal{F}}
\newcommand{\Pbb}{\mathbb{P}}
\newcommand{\Qbb}{\mathbb{Q}}
\newcommand{\Q}{\mathbb{Q}}
\newcommand{\E}{\mathbb{E}}
\newcommand{\Var}{\mathrm{Var}}
\newcommand{\Cov}{\mathrm{Cov}}
\newcommand{\dd}{\,\mathrm{d}}
\newcommand{\1}{\mathbf{1}}
\newcommand{\esssup}{\operatorname*{ess\,sup}}

% Stochastic calc notations
\newcommand{\Ito}{It\^o}
\newcommand{\bracket}[2]{\langle #1,#2\rangle}
\newcommand{\qvar}[1]{\langle #1\rangle}

% ============================================================
% Theorem environments
% ============================================================
\theoremstyle{plain}
\newtheorem{theorem}{Theorem}[section]
\newtheorem{proposition}[theorem]{Proposition}
\newtheorem{lemma}[theorem]{Lemma}
\newtheorem{corollary}[theorem]{Corollary}

\theoremstyle{definition}
\newtheorem{definition}[theorem]{Definition}
\newtheorem{example}[theorem]{Example}
\newtheorem{exercise}[theorem]{Exercise}

\theoremstyle{remark}
\newtheorem{remark}[theorem]{Remark}

% ============================================================
% Minimal lecture header helpers
% ============================================================
\newcommand{\lecture}[2]{%
\section{Lecture #1 --- #2}%
}
\newcommand{\lecturesummary}[1]{%
\begin{center}
\begin{minipage}{0.95\linewidth}
\small\textbf{Summary.} #1
\end{minipage}
\end{center}
\vspace{0.5em}
}

% ============================================================
% Title
% ============================================================
\title{Stochastic Calculus \& Option Pricing --- Synthesis (Weeks 1--7)}
\author{Nathan Sauldubois}
\date{\today}

\begin{document}
\maketitle
\vspace{-0.8em}

\noindent\textbf{How to use these notes (Weeks 1--7).}
\begin{itemize}[leftmargin=1.4em, itemsep=0.2em]
  \item These are \emph{synthetic} lecture notes: definitions, key results, and the computations you actually use.
  \item Each lecture ends with a short checklist: ``what you must be able to do''.
  \item Notation: all processes are defined on a filtered probability space $(\Omega,\F,(\F_t)_{t\ge 0},\Pbb)$ satisfying the usual conditions.
\end{itemize}

\tableofcontents
\bigskip

% ===========================================================

\section{Lecture 1 --- Probability refresher and core mathematical tools}

\lecturesummary{Random variables and their laws; CDFs and densities; moments,
variance, and covariance; dominated convergence; LLN and CLT; characteristic
functions; information and conditional expectation; Gaussian conditioning;
filtrations and martingales; and the differential-calculus formulas used later.}

\paragraph{Reference.}
S.~E. Shreve, \textit{Stochastic Calculus for Finance II}, Chapters~1--2.

% ============================================================
\subsection{Random variables, laws, CDFs, and densities}

\paragraph{Random variable and law.}
On a probability space $(\Omega,\F,\Pbb)$, a real random variable is a
measurable map $X:(\Omega,\F)\to(\R,\mathcal B(\R))$. Its law is the
probability measure
\[
\Pbb_X(A)=\Pbb(X\in A),
\qquad A\in\mathcal B(\R).
\]
The random variable maps states to numbers; its law records the probabilities
of those numbers.

\paragraph{CDF and density.}
The cumulative distribution function of $X$ is
\[
F_X(x)=\Pbb(X\le x).
\]
Every CDF is nondecreasing and right-continuous, with limits $0$ at
$-\infty$ and $1$ at $+\infty$. For $a<b$,
\[
\Pbb(a<X\le b)=F_X(b)-F_X(a).
\]
If $X$ has density $f_X$, then
\[
F_X(x)=\int_{-\infty}^x f_X(u)\dd u,
\qquad
\Pbb(a<X\le b)=\int_a^b f_X(u)\dd u,
\]
and $f_X=F_X'$ wherever the CDF is differentiable. A law and a CDF always
exist; a density need not exist.

\paragraph{Joint and marginal laws.}
The joint law of $(X,Y)$ contains both marginal laws and their dependence.
In the discrete case,
\[
p_X(x)=\sum_y p_{X,Y}(x,y),
\qquad
p_Y(y)=\sum_x p_{X,Y}(x,y),
\]
whereas for a joint density,
\[
f_X(x)=\int_{\R}f_{X,Y}(x,y)\dd y,
\qquad
f_Y(y)=\int_{\R}f_{X,Y}(x,y)\dd x.
\]
Independence is equivalent to factorization of the joint law:
$p_{X,Y}=p_Xp_Y$ or $f_{X,Y}=f_Xf_Y$.

\paragraph{Common one-dimensional laws.}
\begin{center}
\renewcommand{\arraystretch}{1.15}
\begin{tabular}{@{}lccc@{}}
\toprule
Law & Parameters & $\E[X]$ & $\Var(X)$\\
\midrule
Bernoulli & $p$ & $p$ & $p(1-p)$\\
Binomial & $n,p$ & $np$ & $np(1-p)$\\
Poisson & $\lambda$ & $\lambda$ & $\lambda$\\
Uniform & $[a,b]$ & $(a+b)/2$ & $(b-a)^2/12$\\
Exponential & rate $\lambda$ & $1/\lambda$ & $1/\lambda^2$\\
Gaussian & $\mu,\sigma^2$ & $\mu$ & $\sigma^2$\\
\bottomrule
\end{tabular}
\end{center}

\paragraph{Mass functions and densities to recognize.}
For the discrete laws,
\begin{align*}
X\sim\operatorname{Bernoulli}(p):&\quad
\Pbb(X=x)=p^x(1-p)^{1-x},\quad x\in\{0,1\},\\
X\sim\operatorname{Binomial}(n,p):&\quad
\Pbb(X=k)=\binom nkp^k(1-p)^{n-k},\quad k=0,\ldots,n,\\
X\sim\operatorname{Poisson}(\lambda):&\quad
\Pbb(X=k)=e^{-\lambda}\frac{\lambda^k}{k!},\quad k\in\N.
\end{align*}
For the continuous laws,
\begin{align*}
X\sim\operatorname{Unif}[a,b]:&\quad
f_X(x)=\frac1{b-a}\1_{[a,b]}(x),\\
X\sim\operatorname{Exp}(\lambda):&\quad
f_X(x)=\lambda e^{-\lambda x}\1_{[0,\infty)}(x),\\
X\sim\mathcal N(\mu,\sigma^2):&\quad
f_X(x)=\frac1{\sigma\sqrt{2\pi}}
\exp\!\left(-\frac{(x-\mu)^2}{2\sigma^2}\right).
\end{align*}

% ============================================================
\subsection{Expectation, moments, variance, and covariance}

\paragraph{Expectation of a transformation.}
If the relevant sum or integral is finite, then
\[
\E[h(X)]
=\sum_x h(x)\Pbb(X=x)
\quad\text{or}\quad
\E[h(X)]=\int_{\R}h(x)f_X(x)\dd x.
\]
The $k$th moment is $\E[X^k]$. Integrability means
$\E[|X|]<\infty$, and $\E[\1_A]=\Pbb(A)$.

\paragraph{Essential expectation properties.}
For integrable $X,Y$ and constants $a,b$,
\begin{align*}
\E[aX+bY]&=a\E[X]+b\E[Y],\\
X\le Y&\Longrightarrow \E[X]\le\E[Y],\\
|\E[X]|&\le\E[|X|].
\end{align*}
Linearity does not require independence.

\paragraph{Variance and covariance.}
For square-integrable variables,
\[
\Var(X)=\E[(X-\E[X])^2]
=\E[X^2]-\E[X]^2,
\qquad
\operatorname{SD}(X)=\sqrt{\Var(X)},
\]
and
\[
\Cov(X,Y)=\E[XY]-\E[X]\E[Y].
\]
The identities used most often are
\begin{align*}
\Var(aX+b)&=a^2\Var(X),\\
\Var(X+Y)&=\Var(X)+\Var(Y)+2\Cov(X,Y),\\
\Var\!\left(\sum_{k=1}^nX_k\right)
&=\sum_{k=1}^n\Var(X_k)
  +2\sum_{i<j}\Cov(X_i,X_j).
\end{align*}
Independence implies zero covariance, but the converse is false in general.
For independent square-integrable variables, the variance of a sum is the
sum of the variances.

\begin{theorem}[Dominated Convergence Theorem]
Suppose $X_n\to X$ almost surely and there is an integrable random variable
$Y$ such that $|X_n|\le Y$ almost surely for every $n$. Then $X$ is
integrable and
\[
\E[|X_n-X|]\longrightarrow0,
\qquad
\E[X_n]\longrightarrow\E[X].
\]
\end{theorem}

\begin{remark}
Dominated convergence is the standard tool for interchanging a limit and an
expectation. Pointwise convergence alone is not sufficient; the integrable
dominating variable controls the tails uniformly in $n$.
\end{remark}

% ============================================================
\subsection{Sample moments, LLN, and CLT}

For observations $X_1,\ldots,X_n$, define
\[
\overline X_n=\frac1n\sum_{i=1}^nX_i,
\qquad
v_n=\frac1n\sum_{i=1}^n(X_i-\overline X_n)^2,
\qquad
s_n^2=\frac{n}{n-1}v_n.
\]
Here $v_n$ is the empirical variance and $s_n^2$ is its usual unbiased
version under i.i.d. sampling with finite variance.

\begin{theorem}[Law of Large Numbers]
If $(X_i)$ are i.i.d. and $\E[|X_1|]<\infty$, then
\[
\overline X_n\xrightarrow[n\to\infty]{\Pbb}\mu=\E[X_1].
\]
\end{theorem}

\begin{theorem}[Central Limit Theorem]
If, in addition, $0<\sigma^2=\Var(X_1)<\infty$, then
\[
\frac{\sqrt n(\overline X_n-\mu)}{\sigma}
\Rightarrow\mathcal N(0,1).
\]
Thus $\overline X_n$ has approximate standard deviation
$\sigma/\sqrt n$ for large $n$.
\end{theorem}

The LLN says that the estimation error vanishes; the CLT gives its
$n^{-1/2}$ scale and approximate Gaussian shape.

% ============================================================
\subsection{Characteristic functions}

\paragraph{Definition and basic properties.}
The characteristic function of $X$ is
\[
\varphi_X(u)=\E[e^{iuX}],
\qquad u\in\R.
\]
It always exists because $|e^{iuX}|=1$. Moreover,
\[
\varphi_X(0)=1,
\qquad
|\varphi_X(u)|\le1,
\qquad
\varphi_{aX+b}(u)=e^{iub}\varphi_X(au).
\]
The characteristic function uniquely determines the law. If
$\E[|X|^k]<\infty$, then it can be differentiated at the origin and
\[
\varphi_X^{(k)}(0)=i^k\E[X^k].
\]

\paragraph{Independent sums.}
If $X$ and $Y$ are independent, then
\[
\varphi_{X+Y}(u)=\varphi_X(u)\varphi_Y(u).
\]
This turns convolution of laws into multiplication of functions.

\paragraph{Formulas to know.}
\begin{align*}
X\sim\operatorname{Bernoulli}(p):&\quad
\varphi_X(u)=1-p+pe^{iu},\\
X\sim\operatorname{Poisson}(\lambda):&\quad
\varphi_X(u)=\exp\!\bigl(\lambda(e^{iu}-1)\bigr),\\
X\sim\operatorname{Exp}(\lambda):&\quad
\varphi_X(u)=\frac{\lambda}{\lambda-iu},\\
X\sim\mathcal N(\mu,\sigma^2):&\quad
\varphi_X(u)=\exp\!\left(i\mu u-\frac12\sigma^2u^2\right).
\end{align*}

For $Z\sim\mathcal N(0,1)$, differentiation under the integral followed by
integration by parts gives
\[
\varphi_Z'(u)=-u\varphi_Z(u),
\qquad
\varphi_Z(0)=1.
\]
Hence $\varphi_Z(u)=e^{-u^2/2}$; writing $X=\mu+\sigma Z$ gives the general
Gaussian formula.

% ============================================================
\subsection{Information, measurability, and conditional expectation}

\paragraph{$\sigma$-algebras as information.}
A sub-$\sigma$-algebra $\mathcal G\subseteq\F$ represents the events that
can be decided from the available information. A variable $Z$ is
$\mathcal G$-measurable when its value can be determined from that
information. The information revealed by a signal $Y$ is $\sigma(Y)$.

On a finite space, $\mathcal G$ is generated by a partition
$A_1,\ldots,A_m$. A variable is $\mathcal G$-measurable if and only if it is
constant on every atom $A_j$. Equivalently,
\[
Z\text{ is }\sigma(Y)\text{-measurable}
\quad\Longleftrightarrow\quad
Z=g(Y)
\]
for some measurable $g$.

\paragraph{Conditional probability and conditional laws.}
For events $A,B$ with $\Pbb(B)>0$,
\[
\Pbb(A\mid B)=\frac{\Pbb(A\cap B)}{\Pbb(B)}.
\]
For discrete variables and $\Pbb(Y=y)>0$,
\[
\Pbb(X=x\mid Y=y)
=\frac{\Pbb(X=x,Y=y)}{\Pbb(Y=y)}.
\]
For a joint density and $f_Y(y)>0$,
\[
f_{X\mid Y}(x\mid y)
=\frac{f_{X,Y}(x,y)}{f_Y(y)},
\qquad
f_Y(y)=\int_{\R}f_{X,Y}(x,y)\dd x.
\]

\paragraph{Definition of conditional expectation.}
For $X\in L^1$ and $\mathcal G\subseteq\F$, the conditional expectation
$\E[X\mid\mathcal G]$ is the almost-surely unique random variable $Z$ such
that
\[
Z\text{ is }\mathcal G\text{-measurable},
\qquad
\int_AZ\dd\Pbb=\int_AX\dd\Pbb
\quad\text{for every }A\in\mathcal G.
\]
It is the prediction of $X$ based only on the information in $\mathcal G$.
Conditioning on $Y$ means conditioning on $\sigma(Y)$:
\[
\E[X\mid Y]=\E[X\mid\sigma(Y)].
\]

\paragraph{Finite-partition formula.}
If $\mathcal G=\sigma(A_1,\ldots,A_m)$, where the $A_j$ are partition atoms
with positive probability, then
\[
\E[X\mid\mathcal G]
=\sum_{j=1}^m
\frac{\E[X\1_{A_j}]}{\Pbb(A_j)}\1_{A_j}
=\sum_{j=1}^m\E[X\mid A_j]\1_{A_j}.
\]
Conditional expectation is therefore the average of $X$ inside each
information set.

\paragraph{Formula for $\E[h(X)\mid Y]$: discrete case.}
Define, for every $y$ with $\Pbb(Y=y)>0$,
\begin{align*}
m(y)
&=\E[h(X)\mid Y=y]\\
&=\sum_xh(x)\Pbb(X=x\mid Y=y)\\
&=\frac{\sum_xh(x)\Pbb(X=x,Y=y)}{\Pbb(Y=y)}.
\end{align*}
Then
\[
\boxed{\E[h(X)\mid Y]=m(Y).}
\]

\paragraph{Formula for $\E[h(X)\mid Y]$: density case.}
For $f_Y(y)>0$, define
\begin{align*}
m(y)
&=\E[h(X)\mid Y=y]\\
&=\int_{\R}h(x)f_{X\mid Y}(x\mid y)\dd x\\
&=\frac{\int_{\R}h(x)f_{X,Y}(x,y)\dd x}
        {\int_{\R}f_{X,Y}(x,y)\dd x}.
\end{align*}
Again,
\[
\boxed{\E[h(X)\mid Y]=m(Y).}
\]
The final substitution $y=Y$ is essential: a conditional expectation is a
random variable, not merely a number indexed by $y$.

\paragraph{Essential conditional-expectation properties.}
For integrable variables, constants $a,b$, and
$\mathcal H\subseteq\mathcal G\subseteq\F$,
\begin{align*}
\E[aX+bY\mid\mathcal G]
&=a\E[X\mid\mathcal G]+b\E[Y\mid\mathcal G],\\
X\le Y
&\Longrightarrow
\E[X\mid\mathcal G]\le\E[Y\mid\mathcal G],\\
Z\text{ $\mathcal G$-measurable}
&\Longrightarrow
\E[ZX\mid\mathcal G]=Z\E[X\mid\mathcal G],\\
X\text{ $\mathcal G$-measurable}
&\Longrightarrow
\E[X\mid\mathcal G]=X,\\
X\text{ independent of }\mathcal G
&\Longrightarrow
\E[X\mid\mathcal G]=\E[X],\\
\E[\E[X\mid\mathcal G]\mid\mathcal H]
&=\E[X\mid\mathcal H].
\end{align*}
In particular,
\[
\E[\E[X\mid\mathcal G]]=\E[X].
\]
For convex $\phi$, conditional Jensen gives
\[
\phi(\E[X\mid\mathcal G])
\le\E[\phi(X)\mid\mathcal G]
\]
whenever the quantities are integrable.

\paragraph{Gaussian conditioning.}
If $(X,Y)$ is jointly Gaussian, with standard deviations
$\sigma_X,\sigma_Y>0$ and correlation $\rho$, then
\[
X\mid(Y=y)\sim\mathcal N\!\left(
\mu_X+\rho\frac{\sigma_X}{\sigma_Y}(y-\mu_Y),
\;\sigma_X^2(1-\rho^2)
\right).
\]
Consequently,
\[
\E[X\mid Y]
=\mu_X+\rho\frac{\sigma_X}{\sigma_Y}(Y-\mu_Y),
\qquad
\Var(X\mid Y)=\sigma_X^2(1-\rho^2).
\]
The conditional mean is affine in the observed signal, and the residual
variance decreases as $|\rho|$ increases.

% ============================================================
\subsection{Processes, filtrations, and martingales}

\paragraph{Processes and filtrations.}
A discrete-time stochastic process is a sequence $(X_n)_{n\ge0}$ of random
variables on the same probability space. A filtration is an increasing
family of $\sigma$-algebras,
\[
\F_0\subseteq\F_1\subseteq\F_2\subseteq\cdots,
\]
representing information accumulated through time. The natural filtration
of $X$ is
\[
\F_n^X=\sigma(X_0,\ldots,X_n).
\]

A process is adapted if $X_n$ is $\F_n$-measurable for every $n$.
Adaptedness is the no-look-ahead requirement.

\paragraph{Random walks.}
If
\[
S_n=S_0+\sum_{k=1}^n\xi_k,
\]
where the increments are i.i.d. with mean $m$ and variance $\sigma^2$, then
\[
\E[S_n]=S_0+nm,
\qquad
\Var(S_n)=n\sigma^2.
\]

\paragraph{Martingales.}
An integrable adapted process $(M_n)$ is a martingale if
\[
\E[M_{n+1}\mid\F_n]=M_n
\qquad\text{for every }n.
\]
For the random walk above,
\[
M_n=S_n-nm
\]
is a martingale. For independent symmetric $\pm1$ increments,
$S_n^2-n$ is also a martingale. A martingale combines three requirements:
integrability, adaptation, and the conditional fair-game identity.

% ============================================================
\subsection{Differential-calculus review}

\paragraph{One-variable derivative and Taylor expansion.}
For $g:\R\to\R$,
\[
g'(x)=\lim_{h\to0}\frac{g(x+h)-g(x)}{h}.
\]
If $g$ is twice differentiable near $x$, then
\[
g(x+h)=g(x)+g'(x)h+\frac12g''(x)h^2+o(h^2).
\]

\paragraph{Gradient and Hessian.}
For $f:\R^d\to\R$,
\[
\nabla f(x)=
\begin{pmatrix}
\partial_1f(x)\\ \vdots\\ \partial_df(x)
\end{pmatrix},
\qquad
\nabla^2f(x)=\bigl(\partial_{ij}f(x)\bigr)_{1\le i,j\le d}.
\]
The gradient gives first-order sensitivities; the Hessian records curvature
and interactions. The second-order multivariable expansion is
\[
f(x+h)=f(x)+\nabla f(x)^\top h
+\frac12h^\top\nabla^2f(x)h+o(\|h\|^2).
\]

\paragraph{Multivariable chain rule.}
If $x(t)\in\R^d$ and $f$ are differentiable, then
\[
\frac{\dd}{\dd t}f(x(t))
=\nabla f(x(t))^\top x'(t).
\]
In particular, for $V=V(t,s)$ and a differentiable path $S(t)$,
\[
\frac{\dd}{\dd t}V(t,S(t))
=V_t(t,S(t))+V_s(t,S(t))S'(t).
\]
The second-order Taylor term involving $V_{ss}$ will be essential once the
underlying path is Brownian rather than differentiable.

% ============================================================
\subsection*{Lecture 1 checklist}

You should be able to:
\begin{itemize}[leftmargin=1.5em,itemsep=0.15em]
  \item pass between a law, a CDF, a density, and interval probabilities;
  \item compute $\E[h(X)]$, moments, variances, covariances, and variances of sums;
  \item state when dominated convergence allows a limit to pass through $\E$;
  \item distinguish the conclusions of the LLN and the CLT;
  \item compute and use characteristic functions, especially for independent sums;
  \item identify the information in $\sigma(Y)$ and test measurability on partition atoms;
  \item compute $\E[h(X)\mid Y]$ in both discrete and density models;
  \item use linearity, known factors, independence, monotonicity, Jensen, and the tower rule;
  \item apply the Gaussian conditioning formula;
  \item verify adaptation and the martingale property;
  \item compute derivatives, gradients, Hessians, Taylor expansions, and chain rules.
\end{itemize}

% ============================================================

\section{Lecture 2 --- Brownian Motion, Quadratic Variation, and GBM}

\paragraph{References}
% ============================================================

\begin{itemize}[leftmargin=1.4em]

\item \textit{Stochastic Calculus for Finance II: Continuous-Time Models}, S. E. Shreve, Chapter~3, pp.~83-117

\item \textit{Brownian Motion and Stochastic Calculus}, I. Karatzas and S. E. Shreve, Chapter~2, pp.~47-126

\end{itemize}


\subsection{From Discrete Random Walks to Brownian Motion}

\paragraph{Simple symmetric random walk.}
Let $(\varepsilon_k)_{k\ge1}$ be i.i.d.\ with
$\Pbb(\varepsilon_k=\pm1)=\tfrac12$, and define
\[
S_0=0,
\qquad
S_n=\sum_{k=1}^n \varepsilon_k .
\]
Then $(S_n)$ is a discrete-time martingale with independent increments,
$\E[S_n]=0$ and $\Var(S_n)=n$.

\paragraph{Diffusive scaling.}
Introduce a time step $1/n$ and define the rescaled process
\[
W^{(n)}(t)=\frac{1}{\sqrt{n}}\,S_{\lfloor nt\rfloor},
\qquad t\ge0,
\]
with linear interpolation.
By the CLT, increments of $W^{(n)}$ over intervals of length $t-s$
are approximately Gaussian with variance $t-s$.

\paragraph{Donsker invariance principle (informal).}
As $n\to\infty$, $W^{(n)}$ converges in distribution to a continuous-time
process $W$, called \emph{Brownian motion}.
This result explains why Brownian motion is the universal
continuous-time limit of random walks.

\subsection{Definition and Basic Properties of Brownian Motion}

\paragraph{Definition.}
A process $(W_t)_{t\ge0}$ is a \emph{standard Brownian motion} if:
\begin{itemize}
  \item $W_0=0$,
  \item it has independent increments,
  \item for $0\le s<t$, $W_t-W_s\sim\mathcal N(0,t-s)$,
  \item paths are continuous almost surely.
\end{itemize}

\paragraph{Gaussian process viewpoint.}
Brownian motion is a centered Gaussian process fully characterized by
\[
\Cov(W_s,W_t)=\min(s,t).
\]
In particular, for $s<t$,
\[
W_t=W_s+(W_t-W_s),
\]
where $W_s$ and $W_t-W_s$ are independent.

\paragraph{Martingale property.}
With respect to its natural filtration $\mathcal F_t=\sigma(W_s:s\le t)$,
Brownian motion is a martingale:
\[
\E[W_t\mid\mathcal F_s]=W_s,\qquad s\le t.
\]

\subsection{Markov and Strong Markov Properties}

\paragraph{Markov property.}
Brownian motion is Markov: conditionally on $W_s$,
the future is independent of the past,  if $t \geq s$, 
\[
\Pbb(W_t\in A\mid\mathcal F_s)=\Pbb(W_t\in A\mid W_s).
\]
Equivalently,
\[
W_t\mid W_s=x \sim \mathcal N(x,t-s).
\]

\paragraph{Strong Markov property.}
Let $\tau$ be a stopping time.
Conditionally on $\mathcal F_\tau$, the shifted process
\[
(W_{\tau+t}-W_\tau)_{t\ge0}
\]
is a Brownian motion independent of $\mathcal F_\tau$.
This is crucial for hitting times, barriers, and American-style features.

\subsection{Reflection Principle and Hitting Times}

\paragraph{Running maximum.}
Define the running maximum
\[
M_t=\sup_{0\le s\le t} W_s,
\]
and the hitting time of level $a$:
\[
\tau_a=\inf\{t\ge0:\ W_t=a\}.
\]

\paragraph{Reflection principle.}
For $a>0$ and $t>0$,
\[
\Pbb(M_t\ge a)=2\,\Pbb(W_t\ge a).
\]
More generally, for $w\le a$,
\[
\Pbb(M_t>a,\ W_t<w)=\Pbb(W_t>2a-w).
\]

\paragraph{Interpretation.}
Among paths that hit level $a$ before time $t$,
reflecting the path after the first hitting time preserves Wiener measure.
This converts path-dependent events into Gaussian tail probabilities.

\paragraph{Consequence.}
Brownian motion hits every level almost surely:
\[
\Pbb(\tau_a<\infty)=1,
\qquad a\in\R.
\]

\subsection{Pathwise Regularity and Irregularity}

\paragraph{Nowhere differentiability.}
Almost surely, Brownian paths are nowhere differentiable.

\paragraph{Oscillations.}
For every $t\ge0$, almost surely,
\[
\limsup_{h\downarrow0}\frac{W_{t+h}-W_t}{\sqrt{h}}=+\infty,
\qquad
\liminf_{h\downarrow0}\frac{W_{t+h}-W_t}{\sqrt{h}}=-\infty.
\]

\paragraph{Hölder regularity.}
For every $\alpha<\tfrac12$, Brownian motion is locally $\alpha$-Hölder,
but it is not $\tfrac12$-Hölder.

\subsection{Quadratic Variation}

\paragraph{Definition.}
For a partition $\pi=\{0=t_0<\cdots<t_n=T\}$, define
\[
[X]^\pi_T=\sum_{k=0}^{n-1}(X_{t_{k+1}}-X_{t_k})^2.
\]
If this converges as $|\pi|\to0$, the limit is the quadratic variation $[X]_T$.

\paragraph{Finite variation vs.\ Brownian motion.}
If $X_t=\int_0^t v_s ds$, then $[X]_T=0$.
In contrast, for Brownian motion,
\[
[W]_T=T.
\]

\paragraph{Key martingale.}
Quadratic variation explains why
\[
W_t^2-t
\]
is a martingale.
This is the simplest manifestation of Itô calculus.

\paragraph{Quadratic covariation.}
For processes $X,Y$,
\[
[X,Y]_T=\lim_{\pi}\sum_k (X_{t_{k+1}}-X_{t_k})(Y_{t_{k+1}}-Y_{t_k}).
\]
If $A$ has finite variation, then $[A,W]=0$.

\subsection{Heuristic It\^o Rules}

Brownian increments satisfy the scaling:
\[
\Delta W = O(\sqrt{\Delta t}).
\]
As a bookkeeping device,
\[
(dW_t)^2=dt,
\qquad
dW_t\,dt=0,
\qquad
(dt)^2=0.
\]
These rules encode which discrete sums survive in the limit.

\subsection{Geometric Brownian Motion (GBM)}

\paragraph{Motivation.}
Additive models allow negative prices and ignore scale effects.
Instead, model relative changes:
\[
\frac{dS_t}{S_t}=\mu\,dt+\sigma\,dW_t.
\]

\paragraph{Explicit solution.}
The solution is
\[
S_t=S_0\exp\!\Big(
(\mu-\tfrac12\sigma^2)t+\sigma W_t
\Big),
\]
so $S_t$ is lognormally distributed.

\paragraph{Moments.}
\[
\E[S_t]=S_0 e^{\mu t},
\qquad
\Var(S_t)=S_0^2 e^{2\mu t}(e^{\sigma^2 t}-1).
\]

\paragraph{Volatility scaling.}
Log-returns over $\Delta t$ have variance $\sigma^2\Delta t$,
which justifies the $\sqrt{\Delta t}$ rule.

\subsection{Stopping Times and Optional Stopping}

\paragraph{Stopped process.}
For a stopping time $\tau$ and a martingale $(M_t)$, define
\[
M_t^\tau=M_{t\wedge\tau}.
\]
Then $(M_t^\tau)$ is again a martingale.

\paragraph{Optional stopping (safe version).}
If $\tau$ is bounded, then
\[
\E[M_\tau]=\E[M_0].
\]
In general, $\tau<\infty$ a.s.\ is not sufficient;
integrability or uniform integrability is required.

\paragraph{Example.}
For Brownian motion and $\tau_a=\inf\{t:\ W_t=a\}$,
$\E[\tau_a]=\infty$ and $W_{\tau_a}=a$,
so $\E[W_{\tau_a}]\neq\E[W_0]$.
Optional stopping applies only to $t\wedge\tau_a$.

\medskip
\subsection{Multidimensional Brownian Motion}

\paragraph{Definition.}
A $d$-dimensional Brownian motion is a vector-valued process
\[
W_t=(W_t^1,\dots,W_t^d),\qquad t\ge 0,
\]
such that:
\begin{itemize}
  \item $W_0=0\in\R^d$,
  \item it has independent increments,
  \item for $0\le s<t$, the increment $W_t-W_s$ is multivariate Gaussian
  with mean $0$ and covariance matrix $(t-s)I_d$,
  \item paths are continuous almost surely.
\end{itemize}

Equivalently, the components $(W^1,\dots,W^d)$ are independent
one-dimensional Brownian motions.

\paragraph{Gaussian characterization.}
Multidimensional Brownian motion is a centered Gaussian process with
\[
\Cov(W_s^i,W_t^j)=\delta_{ij}\min(s,t),
\qquad 1\le i,j\le d,
\]
where $\delta_{ij}$ is the Kronecker symbol.

\paragraph{Independent vs.\ correlated Brownian motions.}
More generally, one may consider a Brownian motion with covariance matrix
$\Sigma\in\R^{d\times d}$ (symmetric, positive semidefinite), defined by
\[
\Cov(W_s,W_t)=\min(s,t)\,\Sigma.
\]
In this case,
\[
\dd[W^i,W^j]_t=\Sigma_{ij}\,\dd t.
\]
Any such process can be written as
\[
W_t = B\,\widetilde W_t,
\]
where $\widetilde W$ is a standard $d$-dimensional Brownian motion
and $B$ satisfies $BB^\top=\Sigma$ (e.g.\ Cholesky factorization).

\paragraph{Quadratic variation and covariation.}
For a $d$-dimensional Brownian motion $W$,
\[
[W^i]_t = t,
\qquad
[W^i,W^j]_t = 0 \quad (i\neq j)
\]
in the independent case, and more generally
\[
[W]_t = t\,\Sigma,
\]
where $[W]_t$ denotes the matrix-valued quadratic variation.

\paragraph{It\^o rules (multidimensional mnemonic).}
For independent components,
\[
(\dd W_t^i)^2=\dd t,
\qquad
\dd W_t^i\,\dd W_t^j=0 \ (i\neq j),
\qquad
\dd W_t^i\,\dd t=0,
\qquad
(\dd t)^2=0.
\]
With correlation matrix $\Sigma$, this becomes
\[
\dd W_t^i\,\dd W_t^j=\Sigma_{ij}\,\dd t.
\]

\paragraph{Preview.}
In higher dimension, stochastic calculus keeps the same structure:
\emph{drift terms}, \emph{martingale terms}, and \emph{quadratic variation}.
Only the bookkeeping becomes matrix-valued.
This naturally leads to the multidimensional It\^o formula
and to multidimensional diffusion models.

\subsection*{Lecture 2 checklist}

You should be able to:
\begin{itemize}[leftmargin=1.5em,itemsep=0.15em]
  \item state the defining properties of Brownian motion and compute its
        Gaussian increment laws and covariance function;
  \item use independent increments to verify the martingale and Markov
        properties;
  \item apply the reflection principle to convert a maximum or hitting-time
        event into a Gaussian tail probability;
  \item distinguish Brownian paths from finite-variation paths, in particular
        through their quadratic variation;
  \item use the It\^o multiplication rules
        $(\dd W_t)^2=\dd t$, $\dd W_t\dd t=0$, and $(\dd t)^2=0$;
  \item solve the geometric Brownian-motion SDE and identify the
        lognormal law of $S_t$;
  \item identify a stopping time and state a safe version of optional
        stopping, including its hypotheses;
  \item work with independent and correlated multidimensional Brownian
        motions and their covariations.
\end{itemize}

% ============================================================

\section{Lecture 3 --- The It\^o Stochastic Integral}

\lecturesummary{The It\^o integral is the $L^2$ limit of left-point,
non-anticipative gains. Predictability prevents the use of future
information and finite expected energy gives integrability. The resulting
gain process is a square-integrable martingale, with quadratic variation
$\int_0^t\phi_s^2\dd s$; for deterministic integrands, $L^2$ inner products
become covariances of jointly Gaussian variables.}

\subsection{Definition, admissibility, and construction}

\paragraph{Why left endpoints are essential.}
For a grid \(0=t_0<\cdots<t_n=T\), a discrete gain has the form
\[
G_T^\pi=\sum_{i=0}^{n-1}\phi_{t_i}
        (W_{t_{i+1}}-W_{t_i}).
\]
The coefficient used on \((t_i,t_{i+1}]\) must be known at time \(t_i\):
\(\phi_{t_i}\) is \(\F_{t_i}\)-measurable. This is the mathematical content
of predictability and the economic content of non-anticipation. The sampling
rule matters because Brownian quadratic variation survives:
\[
\sum_i W_{t_{i+1}}\Delta W_i-
\sum_i W_{t_i}\Delta W_i
=\sum_i(\Delta W_i)^2\longrightarrow T.
\]

\begin{definition}[Simple predictable integrand]
A simple predictable process on \([0,T]\) is
\[
\phi_t=\sum_{i=0}^{n-1}\xi_i\1_{(t_i,t_{i+1}]}(t),
\qquad \xi_i\in\F_{t_i},\qquad \E[\xi_i^2]<\infty.
\]
Its It\^o integral is the left-point gain
\[
\boxed{
\int_0^T\phi_t\dd W_t
=\sum_{i=0}^{n-1}\xi_i(W_{t_{i+1}}-W_{t_i}).}
\]
\end{definition}

\paragraph{Admissible class and extension.}
The standard \(L^2\) theory assumes
\[
\phi\text{ predictable},
\qquad
\boxed{\E\int_0^T\phi_t^2\dd t<\infty.}
\]
Choose simple predictable \(\phi^n\) such that
\(\E\int_0^T|\phi_t^n-\phi_t|^2\dd t\to0\). Then
\[
\boxed{
\int_0^T\phi_t\dd W_t
:=L^2\text{-}\lim_{n\to\infty}
\int_0^T\phi_t^n\dd W_t.}
\]
The It\^o isometry below proves that the limit exists, is unique, and does
not depend on the approximating sequence.

\subsection{The identities that do the calculations}

Let \(\phi,\psi\) be predictable with finite expected energy. Then
\begin{align*}
\int_0^T(a\phi_t+b\psi_t)\dd W_t
&=a\int_0^T\phi_t\dd W_t+b\int_0^T\psi_t\dd W_t,\\
\E\left[\int_0^T\phi_t\dd W_t\right]&=0,\\
\E\left[\left(\int_0^T\phi_t\dd W_t\right)^2\right]
&=\E\int_0^T\phi_t^2\dd t,\\
\E\left[\left(\int_0^T\phi_t\dd W_t\right)
           \left(\int_0^T\psi_t\dd W_t\right)\right]
&=\E\int_0^T\phi_t\psi_t\dd t.
\end{align*}
The third line is the \emph{It\^o isometry}; the fourth is its polarization
or covariance identity. Thus integrand energy becomes terminal variance.

\begin{theorem}[Martingale gain process]
If \(\phi\) is predictable and
\(\E\int_0^T\phi_s^2\dd s<\infty\), then
\[
M_t=\int_0^t\phi_s\dd W_s
\]
is a continuous square-integrable martingale. In particular, for \(s\leq t\),
\[
\boxed{
\E[M_t\mid\F_s]=M_s,
\qquad
\E\left[\int_s^t\phi_u\dd W_u\middle|\F_s\right]=0.}
\]
\end{theorem}

The quadratic variation and mutual covariation are
\[
\boxed{
[M]_t=\int_0^t\phi_s^2\dd s,
\qquad
\left[\int_0^\cdot\phi_s\dd W_s,
       \int_0^\cdot\psi_s\dd W_s\right]_t
=\int_0^t\phi_s\psi_s\dd s.}
\]
The first formula is pathwise; the isometry is its expected-value analogue.

\begin{remark}[True versus local martingale]
The expected condition
\(\E\int_0^T\phi_s^2\dd s<\infty\) gives a true \(L^2\)-martingale.
Pathwise local square integrability generally gives only a local martingale;
zero drift alone does not justify erasing the stochastic integral after
taking expectations.
\end{remark}

\subsection{Deterministic integrands and formulas to know}

If \(f\in L^2(0,T)\) is deterministic, then
\[
\boxed{
\int_0^Tf(t)\dd W_t
\sim\mathcal N\left(0,\int_0^Tf(t)^2\dd t\right).}
\]
For deterministic \(f,g\), the two integrals are jointly Gaussian and
\[
\Cov\left(\int_0^Tf\dd W,\int_0^Tg\dd W\right)
=\int_0^Tf(t)g(t)\dd t.
\]
Consequently, orthogonality in \(L^2(0,T)\) implies independence of the two
Gaussian integrals. This implication is special to deterministic integrands;
a random predictable integrand need not produce a Gaussian integral.

Three reusable identities are
\[
\boxed{
\int_0^Tc\dd W_t=cW_T,
\qquad
\int_0^Tt\dd W_t\sim\mathcal N\left(0,\frac{T^3}{3}\right),
\qquad
\int_0^TW_t\dd W_t=\frac12(W_T^2-T).}
\]
For deterministic absolutely continuous \(f\), stochastic integration by
parts reduces to
\[
\boxed{
\int_0^Tf(t)\dd W_t
=f(T)W_T-\int_0^TW_tf'(t)\dd t.}
\]
There is no covariation correction because a deterministic finite-variation
function has zero covariation with Brownian motion.

If \(\tau\) is a stopping time, then on every bounded horizon
\[
\boxed{
\int_0^t\1_{\{s\leq\tau\}}\dd W_s=W_{t\wedge\tau},
\qquad
\E[W_{t\wedge\tau}^2]=\E[t\wedge\tau].}
\]

\paragraph{Reusable calculation pattern.}
Before computing an integral, check in order:
\begin{enumerate}[leftmargin=1.6em,itemsep=0.15em]
  \item predictability: is the coefficient known before the increment?
  \item finite expected energy: is
        \(\E\int_0^T\phi_t^2\dd t<\infty\)?
  \item structure: is the integrand deterministic, stopped, or simple?
  \item target: use the isometry for second moments, covariance for joint
        risks, and the martingale property for conditional means.
\end{enumerate}

\subsection*{Lecture 3 checklist}

You should be able to:
\begin{itemize}[leftmargin=1.5em,itemsep=0.15em]
  \item define an It\^o integral first for a simple predictable process;
  \item check predictability and the condition
        $\E\int_0^T\phi_t^2\dd t<\infty$;
  \item use linearity, zero mean, the It\^o isometry, and the covariance
        identity;
  \item explain why $\int_0^t\phi_s\dd W_s$ is a martingale;
  \item identify the normal law of an integral with deterministic integrand;
  \item compute brackets of It\^o integrals;
  \item use deterministic integration by parts and derive
        $\int_0^T W_t\dd W_t=\tfrac12(W_T^2-T)$;
  \item simulate an It\^o integral with left-point Brownian increments and
        validate it with the isometry.
\end{itemize}

% ============================================================

\section{Lecture 4 --- It\^o Processes and It\^o's Formula}

\lecturesummary{Quadratic variation makes stochastic calculus second order,
and covariation turns correlation into a $\dd t$ term. It\^o's formula is a
chain rule plus a curvature correction. Its product and ratio forms solve
SDEs, support discounting and delta hedging, and distinguish local from true
martingales.}

\paragraph{Differential versus integral form.}
For an It\^o process, the compact notation and its precise meaning are
\[
\boxed{
\dd X_t=a_t\dd t+b_t\dd W_t
\quad\Longleftrightarrow\quad
X_t-X_s=\int_s^t a_u\dd u+\int_s^t b_u\dd W_u,
\qquad 0\leq s\leq t.}
\]
Use the differential form for calculations and the integral form as the
rigorous equality. Never divide through by \(\dd t\): Brownian paths are not
differentiable.

\subsection{Quadratic covariation and the differential table}

For continuous processes \(X,Y\), quadratic covariation is the limiting sum
\[
[X,Y]_t=\lim_{|\pi|\to0}
\sum_i(X_{t_{i+1}}-X_{t_i})(Y_{t_{i+1}}-Y_{t_i}),
\]
whenever the limit exists in probability. For Brownian motions \(W,B\) with
instantaneous correlation \(\rho\),
\[
\boxed{
(\dd W_t)^2=\dd t,
\quad \dd W_t\dd B_t=\rho\dd t,
\quad \dd t\dd W_t=0,
\quad (\dd t)^2=0.}
\]
These are mnemonics for quadratic-variation limits, not ordinary algebraic
infinitesimals. If
\[
\dd X_t=a_t\dd t+b_t\dd W_t+c_t\dd B_t,
\quad
\dd Y_t=p_t\dd t+q_t\dd W_t+r_t\dd B_t,
\]
then
\[
\boxed{
\dd[X,Y]_t
=\bigl(b_tq_t+c_tr_t+\rho b_tr_t+\rho c_tq_t\bigr)\dd t.}
\]
Only diffusion coefficients enter; finite-variation drifts contribute no
quadratic covariation.

For the It\^o process
\(\dd X_t=a_t\dd t+b_t\dd W_t\),
\[
\boxed{\dd[X]_t=b_t^2\dd t,
\qquad [X]_t=\int_0^tb_s^2\dd s.}
\]

\subsection{One- and two-dimensional It\^o formulas}

\begin{theorem}[One-dimensional It\^o formula]
Let
\(\dd X_t=a_t\dd t+b_t\dd W_t\) and
\(f\in C^{1,2}\). Then
\[
\boxed{
\dd f(t,X_t)
=\left(f_t+a_tf_x+\frac12b_t^2f_{xx}\right)(t,X_t)\dd t
+b_tf_x(t,X_t)\dd W_t.}
\]
Equivalently,
\(\dd f=f_t\dd t+f_x\dd X+\tfrac12f_{xx}\dd[X]\).
\end{theorem}

The drift vanishes when
\[
\boxed{f_t+a_tf_x+\tfrac12b_t^2f_{xx}=0}
\]
along the process, making \(f(t,X_t)\) a local martingale. A separate
boundedness or integrability argument is required to obtain a true
martingale and use conditional expectations.

\begin{theorem}[Two-dimensional It\^o formula]
Let \(X,Y\) be continuous It\^o processes and let
\(f\in C^{1,2}([0,T]\times\R^2)\), with one continuous time derivative
and continuous spatial derivatives up to order two. Then
\[
\boxed{
\begin{aligned}
\dd f(t,X_t,Y_t)={}&f_t\dd t+f_x\dd X_t+f_y\dd Y_t\\
&+\frac12f_{xx}\dd[X]_t+f_{xy}\dd[X,Y]_t
+\frac12f_{yy}\dd[Y]_t.
\end{aligned}}
\]
All derivatives are evaluated at \((t,X_t,Y_t)\).
In particular, if
\[
\dd X_t=a_t\dd t+b_t\dd W_t,
\qquad \dd Y_t=c_t\dd t+d_t\dd B_t,
\qquad \dd W_t\dd B_t=\rho\dd t,
\]
then
\[
\boxed{
\begin{aligned}
\dd f(t,X_t,Y_t)
={}&\left(f_t+a_tf_x+c_tf_y+\tfrac12b_t^2f_{xx}
       +\rho b_td_tf_{xy}+\tfrac12d_t^2f_{yy}\right)\dd t\\
&+b_tf_x\dd W_t+d_tf_y\dd B_t.
\end{aligned}}
\]
The mixed derivative has no factor \(1/2\); its contribution vanishes
for independent Brownian drivers.
\end{theorem}

\paragraph{Reliable application method.}
Write the state variables as placeholders, compute all first derivatives and
all second derivatives in noisy directions, substitute the processes, then
group separately the \(\dd t\) and Brownian terms. Always check the domain
of \(f\), especially for \(\log x\), \(1/x\), and payoff kinks.

\subsection{Reusable transformation rules}

If \(X_t>0\) and
\(\dd X_t/X_t=\alpha_t\dd t+\beta_t\dd W_t\), then
\[
\boxed{
\dd\log X_t
=\left(\alpha_t-\frac12\beta_t^2\right)\dd t
+\beta_t\dd W_t.}
\]
The term \(-\beta_t^2/2\) is the volatility correction. If
\(\dd X_t=a_t\dd t+b_t\dd W_t\) and \(X_t\neq0\), then
\[
\boxed{
\dd\left(\frac1{X_t}\right)
=\left(-\frac{a_t}{X_t^2}+\frac{b_t^2}{X_t^3}\right)\dd t
-\frac{b_t}{X_t^2}\dd W_t.}
\]

For continuous It\^o processes,
\[
\boxed{
\dd(X_tY_t)=X_t\dd Y_t+Y_t\dd X_t+\dd[X,Y]_t.}
\]

\subsection{Canonical SDEs to solve and reuse}

\paragraph{Geometric Brownian motion.}
For
\[
\dd S_t=\mu S_t\dd t+\sigma S_t\dd W_t,
\qquad S_0>0,
\]
apply It\^o to \(\log S_t\):
\[
\boxed{
S_t=S_0\exp\left\{
\left(\mu-\frac12\sigma^2\right)t+\sigma W_t
\right\}.}
\]
Thus \(S_t>0\), \(\log S_t\) is Gaussian, and
\[
\E[S_t]=S_0e^{\mu t},
\qquad
\frac{\dd(S_t^p)}{S_t^p}
=\left[p\mu+\frac12p(p-1)\sigma^2\right]\dd t
+p\sigma\dd W_t.
\]

\paragraph{Ornstein--Uhlenbeck process.}
For
\[
\dd X_t=\kappa(m-X_t)\dd t+\sigma\dd W_t,
\qquad \kappa>0,
\]
the integrating factor \(e^{\kappa t}\) yields, for \(s\leq t\),
\[
\boxed{
X_t=m+(X_s-m)e^{-\kappa(t-s)}
+\sigma\int_s^te^{-\kappa(t-u)}\dd W_u.}
\]
Therefore
\[
\boxed{
\E[X_t\mid\F_s]=m+(X_s-m)e^{-\kappa(t-s)},
\quad
\Var(X_t\mid\F_s)=\frac{\sigma^2}{2\kappa}
\bigl(1-e^{-2\kappa(t-s)}\bigr).}
\]

\paragraph{Exponential martingale and Brownian moments.}
Two particularly useful drift cancellations are
\[
\boxed{
\exp\left(\lambda W_t-\frac12\lambda^2t\right)
\text{ is a martingale},
\qquad W_t^2-t\text{ is a martingale}.}
\]
More generally, applying It\^o's formula to \(W_t^n\) gives the recursion
\[
\frac{\dd}{\dd t}\E[W_t^n]
=\frac{n(n-1)}2\E[W_t^{n-2}],
\]
whenever the stochastic integral is integrable.

\subsection{Several factors and instantaneous hedging}

If the unhedged shock is
\(\sigma_X\dd W_t-h\sigma_Y\dd B_t\), with
\(\dd W_t\dd B_t=\rho\dd t\), its variance rate is
\[
q(h)=\sigma_X^2+h^2\sigma_Y^2-2h\rho\sigma_X\sigma_Y.
\]
The minimum-variance hedge is
\[
\boxed{
h^*=\rho\frac{\sigma_X}{\sigma_Y},
\qquad q(h^*)=\sigma_X^2(1-\rho^2).}
\]
Only the component correlated with the hedging asset can be removed.

\subsection*{Lecture 4 checklist}

You should be able to:
\begin{itemize}[leftmargin=1.5em,itemsep=0.15em]
  \item translate an It\^o process between differential and integral forms;
  \item identify the time and state variables and write every driving SDE;
  \item compute $f_t$, $f_x$, and $f_{xx}$ before applying It\^o's formula;
  \item apply the two-dimensional formula, including $f_y$, $f_{yy}$,
        and the mixed term $f_{xy}\dd[X,Y]_t$;
  \item calculate covariations, including
        $\dd W_t\dd B_t=\rho\dd t$ for correlated drivers;
  \item retain all terms of order $\dd t$ and separate the final drift and
        diffusion terms;
  \item use product, quotient, and integration-factor formulas for It\^o
        processes;
  \item solve standard SDEs such as GBM and Ornstein--Uhlenbeck processes;
  \item check the domain, $C^{1,2}$ regularity, and integrability conditions;
  \item state the probability measure and financing convention in a financial
        application.
\end{itemize}

% Lectures 5--7 follow the current slide decks and course notes.
% Current synthesis for Lectures 5--7. Included by synthesis.tex.

% ============================================================
\section{Lecture 5 --- Self-Financing, Replication, and Black--Scholes}
% ============================================================

\lecturesummary{It\^o's formula identifies the dynamics of a candidate
price. Self-financing specifies how a trading strategy earns its gains.
Matching those gains with a European claim gives the Black--Scholes PDE
and its delta hedge. The Gaussian solution produces the call formula and
its Greeks. The basic model in this lecture has no dividends.}

\subsection{From It\^o's formula to a pricing equation}

For a stock with historical dynamics
\[
\dd S_t=\mu S_t\dd t+\sigma S_t\dd W_t^{\Pbb},
\qquad
\dd B_t=rB_t\dd t,\qquad B_0=1,
\]
It\^o's formula applied directly to \(V_t=v(t,S_t)\) gives
\[
\dd V_t
=\left(v_t+\mu S_tv_s+\tfrac12\sigma^2S_t^2v_{ss}\right)\dd t
+\sigma S_tv_s\dd W_t^{\Pbb}.
\]
The second derivative term comes from \(\dd[S]_t=\sigma^2S_t^2\dd t\).
For path-dependent payoffs, a running integral such as
\(A_t=\int_0^tS_u\dd u\) can be added to the state. This idea returns in
Lecture~7.

\subsection{The self-financing budget}

Let \(\Delta_t\) be stock shares and \(\beta_t\) bank-account units. The
portfolio value is \(X_t=\Delta_tS_t+\beta_tB_t\). If holdings change at
a trading date \(t_i\), the new shares must be paid for from the bank
account:
\[
S_{t_i}(\Delta_i-\Delta_{i-1})
+B_{t_i}(\beta_i-\beta_{i-1})=0.
\]
The holdings chosen before the next price move earn that move. Passing
from these left-point gains to continuous time gives the
self-financing condition
\[
\boxed{
X_t=X_0+\int_0^t\Delta_u\dd S_u+\int_0^t\beta_u\dd B_u,
\qquad
\dd X_t=\Delta_t\dd S_t+r(X_t-\Delta_tS_t)\dd t.}
\]
The value identity \(X=\Delta S+\beta B\) by itself does not impose
self-financing. For continuous finite-variation holdings, the product rule
also gives \(S_t\dd\Delta_t+B_t\dd\beta_t=0\). For more general
semimartingale holdings, a covariation term must be included in that
rebalancing identity; the stochastic-integral gain equation remains the
definition.

\subsection{Replication, the PDE, and verification}

Suppose a European claim pays \(g(S_T)\) and has smooth candidate price
\(v(t,S_t)\), with \(v(T,s)=g(s)\). A replicating portfolio with
\(X_t=v(t,S_t)\) has gain
\[
\dd X_t
=\left[rv+(\mu-r)S_t\Delta_t\right]\dd t
+\sigma S_t\Delta_t\dd W_t^{\Pbb}.
\]
Matching its Brownian coefficient with that of \(\dd v(t,S_t)\) gives
\(\Delta_t=v_s(t,S_t)\). Matching drifts then cancels \(\mu\) and yields
\[
\boxed{
v_t+rs\,v_s+\tfrac12\sigma^2s^2v_{ss}-rv=0,
\qquad v(T,s)=g(s).}
\]
The cash holding is
\[
\boxed{\beta_t=\frac{v(t,S_t)-S_tv_s(t,S_t)}{B_t}.}
\]
Conversely, a sufficiently regular PDE solution with an admissible
self-financing strategy and the correct terminal payoff verifies the
replication. No arbitrage then identifies its initial cost as the claim
price. The argument uses one traded Brownian exposure to span the one
modeled Brownian risk.

\subsection{Gaussian solution, call formula, and Greeks}

Set \(\tau=T-t\). The logarithmic change of variable and removal of the
discounting term transform the PDE into a heat equation. Its Gaussian
representation is
\[
v(t,s)=e^{-r\tau}
\E\!\left[
g\!\left(s e^{(r-\sigma^2/2)\tau+\sigma\sqrt{\tau}Z}\right)
\right],
\qquad Z\sim N(0,1).
\]
For a call with strike \(K>0\), put
\[
d_1=\frac{\log(s/K)+(r+\sigma^2/2)\tau}{\sigma\sqrt{\tau}},
\qquad d_2=d_1-\sigma\sqrt{\tau}.
\]
The call, put, and parity formulas are
\[
\boxed{\begin{aligned}
C(t,s)&=s\Phi(d_1)-Ke^{-r\tau}\Phi(d_2),\\
P(t,s)&=Ke^{-r\tau}\Phi(-d_2)-s\Phi(-d_1),\\
C(t,s)-P(t,s)&=s-Ke^{-r\tau}.
\end{aligned}}
\]
Here \(\Phi\) and \(\varphi=\Phi'\) are the standard normal CDF and
density. The identity
\[
\boxed{s\varphi(d_1)=Ke^{-r\tau}\varphi(d_2)}
\]
cancels the density terms when differentiating. With calendar-time
\(\Theta=\partial_tC\), the call Greeks are
\[
\boxed{\begin{aligned}
\Delta_C&=\Phi(d_1),&
\Gamma_C&=\frac{\varphi(d_1)}{s\sigma\sqrt{\tau}},&
\mathrm{Vega}_C&=s\varphi(d_1)\sqrt{\tau},\\
\Theta_C&=-\frac{s\sigma\varphi(d_1)}{2\sqrt{\tau}}
           -rKe^{-r\tau}\Phi(d_2),&
\rho_C&=K\tau e^{-r\tau}\Phi(d_2).
\end{aligned}}
\]
Delta is the stock holding in the hedge. Gamma measures how quickly
that holding changes. Discrete rebalancing, funding frictions, jumps,
and volatility misspecification make the continuous-time hedge
imperfect.

\subsection*{Lecture 5 checklist}
\begin{itemize}[leftmargin=1.5em,itemsep=0.15em]
  \item distinguish portfolio value from self-financing gains;
  \item match diffusion and drift to obtain delta and the pricing PDE;
  \item recover the call formula and differentiate it to obtain the Greeks;
  \item check regularity, integrability, and the terminal payoff when
        verifying a candidate hedge.
\end{itemize}

% ============================================================
\section{Lecture 6 --- Reweighting, Girsanov, and Risk-Neutral Measure}
% ============================================================

\lecturesummary{A change of measure visits events with different
frequencies while an inverse density preserves the original
expectation. Gaussian tilting becomes Girsanov's drift shift for
Brownian paths. A true density martingale is needed before the
risk-neutral measure can be used.}

\subsection{Rare tails motivate a change of weights}

A change of measure can make a rare event easier to sample. To recover
its original probability, reweight the observations by the inverse density:
\[
\Pbb(A)=\E^{\Qbb}\!\left[
\frac{\dd\Pbb}{\dd\Qbb}\mathbf1_A\right]
\]
for equivalent measures \(\Pbb\) and \(\Qbb\).

\subsection{Densities and conditional Bayes' formula}

If \(Z_T>0\), \(\E^{\Pbb}[Z_T]=1\), and
\(\dd\Qbb/\dd\Pbb|_{\F_T}=Z_T\), then for integrable \(Y\),
\[
\boxed{\E^{\Qbb}[Y]=\E^{\Pbb}[Z_TY].}
\]
The density process is \(Z_t=\E^{\Pbb}[Z_T\mid\F_t]\). For
\(\mathcal G\subseteq\F_T\), the conditional version is
\[
\boxed{
\E^{\Qbb}[Y\mid\mathcal G]
=\frac{\E^{\Pbb}[Z_TY\mid\mathcal G]}
       {\E^{\Pbb}[Z_T\mid\mathcal G]}.}
\]
Changing measure changes the weights of paths, while equivalent
measures agree on null events.

\subsection{Stochastic exponential and the sign of Girsanov}

For a predictable process \(\lambda\) with finite pathwise energy, the
candidate density is
\[
\boxed{
Z_t=\exp\!\left(
-\int_0^t\lambda_u\dd W_u^{\Pbb}
-\tfrac12\int_0^t\lambda_u^2\dd u\right),
\qquad
\dd Z_t=-Z_t\lambda_t\dd W_t^{\Pbb}.}
\]
This is initially a nonnegative local martingale. A sufficient
condition for it to be a true martingale on \([0,T]\) is Novikov's
condition
\[
\E^{\Pbb}\!\left[
\exp\!\left(\tfrac12\int_0^T\lambda_u^2\dd u\right)\right]<\infty.
\]
If the density is a true martingale and
\(\dd\Qbb/\dd\Pbb|_{\F_T}=Z_T\), Girsanov gives
\[
\boxed{
W_t^{\Qbb}=W_t^{\Pbb}+\int_0^t\lambda_u\dd u.}
\]
Thus \(\dd W_t^{\Pbb}=\dd W_t^{\Qbb}-\lambda_t\dd t\). If
\(\dd X_t=b_t\dd t+a_t\dd W_t^{\Pbb}\), then
\[
\dd X_t=(b_t-a_t\lambda_t)\dd t+a_t\dd W_t^{\Qbb}.
\]
The drift changes; the quadratic variation
\(\dd[X]_t=a_t^2\dd t\) does not.

\subsection{The Black--Scholes pricing measure}

For the non-dividend-paying stock in Lecture~5, set
\(\lambda=(\mu-r)/\sigma\). With the corresponding true density
martingale,
\[
\boxed{
\frac{\dd S_t}{S_t}=r\dd t+\sigma\dd W_t^{\Qbb},
\qquad
\dd(e^{-rt}S_t)=\sigma e^{-rt}S_t\dd W_t^{\Qbb}.}
\]
The discounted price is a local martingale, and in this constant
coefficient model it is a true martingale. The measure \(\Qbb\)
is a pricing measure, not a forecast of historical returns. Its
conditional lognormal law also yields
\(C_t=S_t\Phi(d_1)-Ke^{-r(T-t)}\Phi(d_2)\) from two Gaussian
tails.

\subsection*{Lecture 6 checklist}
\begin{itemize}[leftmargin=1.5em,itemsep=0.15em]
  \item write a density and its inverse before reweighting an expectation;
  \item distinguish a local exponential from a true density martingale;
  \item keep the sign of the Brownian drift shift consistent with \(Z\);
  \item derive the risk-neutral drift of a non-dividend-paying stock.
\end{itemize}

% ============================================================
\section{Lecture 7 --- Martingale Prices, Hedges, and Completeness}
% ============================================================

\lecturesummary{Discounted self-financing gains lead to conditional
valuation when the relevant local martingale is a true martingale.
Dividends enter the stock gain and change its risk-neutral price drift.
Changing numeraire changes the units and measure together. Martingale
representation identifies a claim's Brownian exposure; traded assets
must span that exposure to implement a hedge.}

\subsection{From stochastic integrals to martingale prices}

Let \(B_t=\exp(\int_0^tr_u\dd u)\) and \(D_t=B_t^{-1}\).
For a non-dividend-paying stock under a money-market pricing measure,
\[
\frac{\dd S_t}{S_t}=r_t\dd t+\sigma_t\dd W_t^{\Qbb}
\quad\Longrightarrow\quad
\dd(D_tS_t)=D_t\sigma_tS_t\dd W_t^{\Qbb}.
\]
The discounted price is a local martingale because it is a stochastic
integral. For a self-financing portfolio
\(X_t=\Delta_tS_t+\beta_tB_t\), integration by parts gives
\[
\boxed{\dd(D_tX_t)=\Delta_t\dd(D_tS_t).}
\]
If this discounted wealth is a true martingale and \(X_T=H\), then
\[
\boxed{
V_t=X_t=B_t\E^{\Qbb}\!\left[\frac{H}{B_T}\middle|\F_t\right].}
\]
Local martingality alone does not justify the conditional-expectation
equality; integrability or an appropriate admissibility argument is
needed.

\subsection{Dividend gains and the pricing equation}

Suppose the stock pays continuous dividends \(qS_t\dd t\).
Its \emph{historical ex-dividend price} still has the form
\[
\dd S_t=\mu S_t\dd t+\sigma S_t\dd W_t^{\Pbb}.
\]
The gain from one share is \(\dd S_t+qS_t\dd t\). A
self-financing hedge therefore satisfies
\[
\boxed{
\dd X_t=\Delta_t(\dd S_t+qS_t\dd t)+\beta_t\dd B_t,}
\]
and its discounted form is
\[
\dd(D_tX_t)=\Delta_t\left[
\dd(D_tS_t)+D_tqS_t\dd t\right].
\]
For \(V_t=v(t,S_t)\), matching Brownian coefficients gives
\(\Delta_t=v_s\). Matching drifts cancels the historical price
drift \(\mu\) and gives
\[
\boxed{
v_t+(r-q)sv_s+\tfrac12\sigma^2s^2v_{ss}-rv=0.}
\]
The matching Brownian shift is
\(\lambda=(\mu+q-r)/\sigma\), so under the associated pricing measure
\[
\frac{\dd S_t}{S_t}=(r-q)\dd t+\sigma\dd W_t^{\Qbb}.
\]
The discounted \emph{gain} is driftless; the discounted
ex-dividend price alone has dividend-related drift.

\subsection{Change of numeraire}

For a strictly positive traded numeraire \(N\) with associated measure
\(\Qbb^N\), a replicable payoff has the representation
\[
\boxed{
V_t=N_t\E^{\Qbb^N}\!\left[\frac{H}{N_T}\middle|\F_t\right].}
\]
The numeraire, its measure, and the ratio \(H/N_T\) form one matched
choice of units. Taking \(N=B\) recovers money-market discounting.

\subsection{Martingale representation turns prices into hedges}

In a Brownian filtration, a square-integrable discounted claim value
has a representation
\[
M_t=\E^{\Qbb}[D_TH\mid\F_t]
=M_0+\int_0^t\psi_u^\top\dd W_u^{\Qbb}.
\]
The vector \(\psi_t\) is the claim's exposure to Brownian shocks. With
one non-dividend-paying traded stock satisfying
\(\dd(D_tS_t)=D_t\sigma_tS_t\dd W_t^{\Qbb}\), exposure matching gives
\[
\boxed{\Delta_t=\frac{\psi_t}{D_t\sigma_tS_t}.}
\]
If the claim is \(v(t,S_t)\), this reduces to \(\Delta_t=v_s(t,S_t)\).
For \(m\) assets driven by \(d\) Brownian factors, write
\(\dd\bar S_t=\Gamma_t\dd W_t^{\Qbb}\), with
\(\Gamma_t\in\R^{m\times d}\). Holdings \(\phi_t\) implement a
claim exposure \(\psi_t\) when
\[
\boxed{\Gamma_t^\top\phi_t=\psi_t.}
\]

\subsection{Completeness and the spanning condition}

In a \(d\)-factor Brownian market, all Brownian exposures can be spanned
when the discounted volatility matrix satisfies
\[
\boxed{\operatorname{rank}(\Gamma_t)=d
\quad\text{for almost every }(t,\omega).}
\]
Together with the required admissibility and integrability conditions,
this is the completeness criterion. A new traded asset helps only if
its Brownian exposure adds a missing direction.

When a Brownian direction is unspanned, its risk-premium shift can remain
free across equivalent martingale measures. Claims exposed to that
direction need not have a unique no-arbitrage price.

\subsection{Path dependence and the choice of state}

The full-filtration pricing identity applies to path-dependent claims.
To express the price as a function of the current state, include enough
variables to record the relevant path information. Apply the
multidimensional It\^o formula to this enlarged state: a finite-variation
state contributes a first derivative but no quadratic-variation or
covariation term.

\subsection*{Lecture 7 checklist}
\begin{itemize}[leftmargin=1.5em,itemsep=0.15em]
  \item distinguish discounted prices from discounted dividend gains;
  \item derive the dividend-adjusted PDE and delta hedge by replication;
  \item match a numeraire with its measure in a valuation formula;
  \item use martingale representation and a rank test to identify
        hedgeable Brownian exposures;
  \item identify the additional state variables needed for a path-dependent claim.
\end{itemize}

\paragraph{Examples to know.}
\begin{itemize}[leftmargin=1.5em,itemsep=0.15em]
  \item Price a unit payoff, a terminal share, a forward, and a digital;
        recover a European call by decomposing its payoff.
  \item Price and hedge a call with continuous dividends; recover
        put--call parity using a prepaid share.
  \item Choose a numeraire that simplifies the terminal payoff ratio.
  \item Find the martingale representations of linear and quadratic
        Brownian terminal payoffs.
  \item Analyze a stock driven by two Brownian risks and a claim depending
        on an untraded sector state; determine whether a second traded
        product completes the market.
  \item Use a nontraded claim to illustrate price ambiguity in an
        incomplete market.
  \item Introduce the running average for an arithmetic Asian option
        and derive its two-state pricing PDE.
\end{itemize}


\end{document}
